% Figure from MATH 452 notes, section 8. Compile with the notes preamble.
\begin{tikzpicture}[scale=1.35,
  declare function={f(\t) = 0.5 + 0.18*\t*\t;}]
  \def\xa{1.2} \def\xb{3.0}
  % axes
  \draw[-stealth, thin] (-0.1,0) -- (3.9,0);
  \draw[-stealth, thin] (0,-0.1) -- (0,2.75);
  % tangent line at x (slope f'(x) = 0.36 * 1.2 = 0.432)
  \draw[black!55, thin, dashed] (0.2,{f(\xa)-0.432}) -- (3.55,{f(\xa)+0.432*2.35});
  % graph of f
  \draw[figB, thick, domain=0.15:3.45, samples=60, smooth] plot (\x,{f(\x)});
  \node[figB, font=\small] at (3.6,2.55) {$f$};
  % increment h = dx
  \draw[black!55, thin] (\xa,{f(\xa)}) -- (\xb,{f(\xa)}) node[midway, below, font=\scriptsize, black] {$h = dx$};
  % df = f'(x) h : rise along the tangent
  \draw[figR, very thick] (\xb,{f(\xa)}) -- (\xb,{f(\xa)+0.432*1.8});
  \node[figR, font=\scriptsize, anchor=east] at (\xb-0.06,{f(\xa)+0.15}) {$df = f'(x)\,h$};
  % the remainder between tangent and graph
  \draw[black!40, thin, dotted] (\xb,{f(\xa)+0.432*1.8}) -- (\xb,{f(\xb)});
  % f(x+h) - f(x): brace on the right
  \draw[decorate, decoration={brace, amplitude=3pt}, black!70] (\xb+0.1,{f(\xb)}) -- (\xb+0.1,{f(\xa)})
    node[midway, right, font=\scriptsize, xshift=2pt] {$f(x+h)-f(x)$};
  % points
  \filldraw (\xa,{f(\xa)}) circle (1.2pt);
  \filldraw (\xb,{f(\xb)}) circle (1.2pt);
  \filldraw[figR] (\xb,{f(\xa)+0.432*1.8}) circle (1.1pt);
  % ticks
  \draw[black!40, thin, dotted] (\xa,0) -- (\xa,{f(\xa)});
  \draw[black!40, thin, dotted] (\xb,0) -- (\xb,{f(\xa)});
  \node[font=\scriptsize, below] at (\xa,0) {$x$};
  \node[font=\scriptsize, below] at (\xb,0) {$x+h$};
\end{tikzpicture}
